\Needspace{13\baselineskip}
\chapterrule
\section{分类损失：预测错了哪一类，怎样调整分数？}
\label{intro:20c-classification-loss}

预测下一 Token 是一个\textbf{多类别预测}问题：
模型要为词表里的每个候选 ID 打分，据此比较下一 Token 的候选。
回归问题可以用均方误差（MSE）；这里先算适用于类别概率的交叉熵，
再研究 Embedding 和网络怎样产生这些分数。

先说明本节的范围：这里先用一个三分类例子说明“分数、概率和损失”怎样连接。
上一章已规定主线五项词表的输入与目标怎样对齐；
这里先用独立的三分类小例子算清评分过程，再回到五项词表。

\subsection{先给每一类一个分数：Logits}
\phantomsection\label{concept:logits}

假设只有三类，模型输出
\[
\vect z=(0,\ln 2,0)\quad(1\times3).
\]
这三个数称为 \textbf{logits（未归一化分数）}，
是实数分数；经过 Softmax 才得到概率。
输出维度为类别数：三分类有三个 logits；
本书五项词表的下一 Token 预测则有五个 logits。
二分类可以采用两个 logits，也可以采用一个 logit 加 Sigmoid，
两种参数化方式都常见。

\subsection{Softmax：把分数变成一份概率}
\phantomsection\label{concept:softmax}

分数可以为负，也不要求加起来等于 $1$；概率则必须非负、总和为 $1$。
直接除以分数总和可行吗？另取分数 $(2,-1,0)$，总和为 $1$，相除后仍有负数。
我们需要先把每项变成正数，再分配总共为 $1$ 的概率。

回到本章的 $(0,\ln2,0)$。先逐项取指数，得到 $(1,2,1)$：
每项都为正，而且原来分数较大的项仍较大。
再除以总和 $1+2+1=4$，得到 $(1/4,1/2,1/4)$。
这套“取指数，再除以总和”的运算叫作\textbf{Softmax}。
对每个类别 $j$，一般写成
\begin{equation}
\boxed{p_j=\frac{e^{z_j}}{\sum_{k=1}^{C}e^{z_k}}}\qquad
(C=\text{类别数}).
\label{eq:early-class-softmax}
\end{equation}
这种运算对分数可微，后面可以通过它反传梯度。
\Needspace{90mm}
% Show every score becoming a positive weight, then a probability bar.
\tagstructbegin{tag=Figure,alttext={Softmax 的三步对应图：三类分数为零、自然对数二、零；取指数得到一、二、一；各除以总和四，得到零点二五、零点五零、零点二五。三条概率柱长度按一比二比一绘制，总和为一。}}
\begin{center}
\begin{tikzpicture}[x=1mm,y=1mm,>=Latex,
 n/.style={draw=linegray,rounded corners=2pt,fill=paper,minimum width=18mm,minimum height=9mm,font=\small}]
\node[font=\small\bfseries] at (0,13) {分数 $z$};
\node[font=\small\bfseries] at (43,13) {指数 $e^z$};
\node[font=\small\bfseries] at (104,13) {概率 $p$};
\foreach \y/\z/\e/\width/\prob/\class in {0/0/1/15/0.25/1,-16/{\log2}/2/30/0.50/2,-32/0/1/15/0.25/3}{
 \node[font=\footnotesize,anchor=east] at (-13,\y) {类 \class};
 \node[n] (z\class) at (0,\y) {$\z$};
 \node[n,fill=blueback] (e\class) at (43,\y) {$\e$};
 \draw[->,thick,blue] (z\class)--(e\class) node[midway,above,font=\scriptsize] {取指数};
 \draw[->,thick,teal] (e\class.east)--(79,\y) node[midway,above,font=\scriptsize] {除以 $4$};
 \filldraw[fill=purple!65,draw=purple] (82,{\y-3.5}) rectangle ({82+\width},{\y+3.5});
 \node[font=\small,anchor=west] at ({84+\width},\y) {$\prob$};
}
\node[font=\footnotesize,text=muted] at (55,-44) {分母 $1+2+1=4$；最终 $0.25+0.50+0.25=1$};
\end{tikzpicture}
\captionof{figure}{Softmax 逐项对应：分数取指数，再除以共同总和，得到和为一的权重}
\label{fig:softmax-prob-bars}
\end{center}
\tagstructend


Softmax 只在\textbf{类别维度}上归一化；后面 Attention
也会用同一运算把匹配分数变成读取权重，不过那里归一化的是可见 Key，
不是词表候选类别。


\subsection{交叉熵：正确类的概率越小，惩罚越大}
\phantomsection\label{concept:cross-entropy}

得到概率后，怎样评价这次预测？只看概率最大的类别，会漏掉置信程度的变化：
即使选出的类别相同，给正确答案 $0.9$ 与 $0.4$ 的概率，也应得到不同的分数。
我们希望正确答案概率为 $1$ 时不扣分，概率越低扣分越多。
一种办法是用 $-\ln p_{\rm correct}$ 计分：
\[
p_{\rm correct}=1,\ 1/2,\ 1/4
\quad\longrightarrow\quad
0,\ \ln2,\ \ln4.
\]
前面给连续数值预测用过均方误差；这里选择直接评价正确类概率的\textbf{交叉熵}。
选择对数还因为它能把多次预测的代价相加：若连续两次都给正确答案 $1/2$ 的条件概率，
两次都对的概率是 $1/2\times1/2=1/4$，而 $-\ln(1/4)=\ln2+\ln2$。
下面先写单次预测的损失，下一节再把这个性质用于一串 Token。

本例真实标签为第二类，概率为 $1/2$，损失就是 $\ln2\approx0.6931$。
把“第二类正确”写成只有第二项为 $1$、其余为 $0$ 的 One-hot 向量
$\vect y=(0,1,0)$，就能用求和统一表示这项损失：
\begin{equation}
\boxed{L=-\sum_{j=1}^{C}y_j\ln p_j=-\ln p_{\rm correct}.}
\label{eq:early-class-cross-entropy}
\end{equation}
求和中的 $y_j$ 只选中正确类别；这里的 $\ln$ 是自然对数。

\subsection{为什么取对数：把一串概率变成可以相加的惊讶度}
正确答案的概率越小，看到它时越“意外”。用 $-\ln p$ 表示这份惊讶：
必然事件 $p=1$ 的代价为零；概率每缩小一半，代价都增加同一个 $\ln2$。
例如 $p=1/2,1/4,1/8$ 对应 $\ln2,2\ln2,3\ln2$。

对自回归序列，整段真实答案的条件概率是各位置概率的乘积：
\[
p_\theta(y_1,\ldots,y_T)=\prod_{t=1}^T p_\theta(y_t\mid y_{<t}).
\]
最大化这份\textbf{似然}，等价于最小化负对数：
\[
-\ln p_\theta(y_1,\ldots,y_T)
=\sum_{t=1}^T-\ln p_\theta(y_t\mid y_{<t}).
\]
对数把乘积化为和，便于分位置计分和反传，也避免直接连乘许多小数。
本章先用平均负对数损失衡量预测。
第~\ref*{intro:23e-tech-judgement}~章再介绍由它得到的困惑度，以及评测时需要统一的口径。

\begin{example}
\noindent\emph{{\symboliconfont ★} 选读：计算机怎样避免指数溢出？}\par
面对 logits $(1000,1001)$，直接计算 $e^{1000}$ 容易溢出。
令 $m=\max_jz_j=1001$，统一减去 $m$ 后计算 Softmax：
\[
\frac{(e^{-1},1)}{e^{-1}+1}\approx(0.2689,0.7311),\qquad
p_j=\frac{e^{z_j-m}}{\sum_k e^{z_k-m}}.
\]
分子分母同时除以 $e^m$ 不改变概率，每个指数的输入都不大于零，且至少有一项指数为 $1$。
把两项都再加 $5000$，最大值也增加 $5000$，相减后仍是 $(-1,0)$。

若正确类是第一类，实现还可直接从 logits 计算交叉熵：
\[
L=(m-z_{\rm correct})+\ln\sum_k e^{z_k-m}
=1+\ln(1+e^{-1})\approx1.3133.
\]
这种 LogSumExp 写法避免先求极小概率再取对数，减少概率下溢成零的问题。
Attention 的逐行 Softmax 也使用同样的减最大值技巧。
\end{example}

\subsection{\texorpdfstring{为什么从 logits 反传得到 $p-y$？}{为什么从 logits 反传得到 p 减 y？}}
\phantomsection\label{concept:logit-gradient}

先只求正确类别是第二类的这一道题。将 Softmax 的概率代入损失：
\[
L=-\ln\frac{e^{z_2}}{e^{z_1}+e^{z_2}+e^{z_3}}
=-z_2+\ln(e^{z_1}+e^{z_2}+e^{z_3}).
\]
这里用了 $\ln(a/b)=\ln a-\ln b$ 和 $\ln(e^{z_2})=z_2$。
损失分成两部分：单独的 $-z_2$，以及依赖所有分数的对数项。

求导还需要两条规则：$e^z$ 的导数仍是 $e^z$；
$\ln u$ 对 $u$ 的导数是 $1/u$（$u>0$）。
例如对 $z_2$ 求导，对数项要按链式法则先乘外层的倒数，再乘内层的导数：
\[
\frac{\partial}{\partial z_2}\ln(e^{z_1}+e^{z_2}+e^{z_3})
=\frac{1}{e^{z_1}+e^{z_2}+e^{z_3}}\,e^{z_2}=p_2.
\]
回到 $(0,\ln2,0)$，指数值是 $(1,2,1)$，分母为 $4$。
正确类 $z_2$ 还影响单独的 $-z_2$ 项，所以
\[
\frac{\partial L}{\partial z_2}=-1+\frac24=-\frac12.
\]
错误类 $z_1$ 不出现在 $-z_2$ 中，这一部分导数为零，因此
\[
\frac{\partial L}{\partial z_1}=0+\frac14=\frac14;
\qquad \frac{\partial L}{\partial z_3}=\frac14.
\]
一般情况下，One-hot 的 $y_j$ 记录第 $j$ 类是否正确：正确类为 $1$，其余为 $0$。
把刚才两种情况合起来，就得到
\[
\frac{\partial L}{\partial z_j}
=-y_j+\frac{e^{z_j}}{\sum_k e^{z_k}}
=\boxed{p_j-y_j}.
\]
因此三分类例子的梯度是
\[
\nabla_{\vect z}L=(1/4,-1/2,1/4).
\]
正确类的梯度为负，梯度下降会倾向于提高它的 logit；
另外两项为正，会倾向于降低它们。
对 $N$ 个等权位置的\textbf{平均交叉熵}，
每个位置的这项梯度还要除以 $N$。
这正是后面三个 Token 位置联合训练时出现 $1/3$ 的原因。

\subsection{回到主线：五个候选怎样计分？}

沿用上一章已经确定的输入与目标：
\[
[2,4,5]\longrightarrow[4,5,3].
\]
主线的五项教学词表就是 $C=|\mathcal V|=5$ 个类别，
每个输入位置输出一行五个 logits，逐行 Softmax 后排成概率矩阵 $\Pi$，
第 $i$ 行第 $j$ 列记作 $\Pi_{i,j}$。
读到“春天”时，下一 Token 应是“花”，即第一行的目标 ID 为 $4$
（三个位置的完整对齐见\hyperref[intro:20b-task-contract]{预测任务章}）。
交叉熵只读取目标列：
\[
L_1=-\ln\Pi_{1,4}.
\]
若五个 logits 相等，$\Pi_{1,4}=1/5$，$L_1=\ln5\approx1.609$；
若另一次预测给目标分配 $1/2$ 的概率，则 $L_1=\ln2\approx0.693$，损失更小。
本章三分类例的 $p-y$ 梯度照样适用：只是 $\vect y$ 换成第 $4$ 列为 $1$ 的五维 one-hot。

\begin{warningbox}
“输出维度等于类别数”适用于这里的多类分类输出头，
不是所有任务都必须改写成分类；连续数值预测等任务可以使用其他输出与损失。
\end{warningbox}

\chapterchecktasks{20c-classification-loss}{理解与复算}{
\item \textbf{辨析：}三类 logits 为 $(2,-1,0)$。为什么它们不是合法概率，Softmax 要满足哪些条件？\chapterinlinehint{核对非负性与归一化，再对照原始分数。}
\item \textbf{手算：}三分类概率 $(1/4,1/2,1/4)$，正确类为第二类。求交叉熵以及正确类 logit 的梯度。\chapterinlinehint{使用负对数损失与正确类的 $p-y$。}
\item \textbf{预测：}标签固定不变时，正确类概率从 $0.2$ 升到 $0.4$，交叉熵的变化方向是什么？\chapterinlinehint{比较负对数函数在这两个概率处的大小。}
}


\par\smallskip
分类任务先输出 logits，再用 Softmax 得到概率，
用交叉熵衡量正确类概率的高低；logit 梯度为 $p-y$。
接下来要问：Token ID 如何变成神经网络可学习的输入？
